\subsection*{JuMP}

JuMP (Julia for Mathematical Programming) é um ambiente de modelagem de código aberto para otimização matemática em Julia. Ele permite que os usuários formulem problemas de otimização de forma expressiva e de alto nível, com suporte a diferentes tipos de problemas (programação linear, inteira, quadrática, não linear, etc.) e a uma vasta gama de solucionadores (solvers) comerciais e open-source. JuMP se destaca por sua flexibilidade, performance e pela integração com a linguagem Julia, que combina a facilidade de uso de linguagens de script com a velocidade de linguagens compiladas. A capacidade de estender o JuMP com pacotes adicionais, como o HiGHS para problemas lineares e inteiros, o torna uma ferramenta poderosa para pesquisa e aplicações práticas em otimização (Lubin et al., 2023).

\subsection*{SDDP.jl}

SDDP.jl é um pacote em Julia que implementa a Programação Dinâmica Dual Estocástica (PDDE) de forma eficiente e flexível. Desenvolvido sobre o JuMP, ele facilita a modelagem e a resolução de problemas de programação estocástica multiestágio que possuem a estrutura de um grafo de política. SDDP.jl abstrai muitos dos detalhes complexos da implementação da PDDE, como a geração e o gerenciamento de cortes de Benders, permitindo que os usuários se concentrem na formulação do problema. Ele oferece suporte a uma ampla gama de funcionalidades, incluindo a capacidade de lidar com diferentes distribuições de incerteza e a integração com vários solucionadores. A utilização de SDDP.jl tem sido fundamental para avançar a pesquisa e a aplicação da PDDE em problemas reais, como o planejamento de sistemas de energia e gestão de portfólio (Dowson e Kapelevich, 2021).

\subsection*{HiGHS}

HiGHS (High-performance Gurobi-like Heuristic Solver) é um resolvedor (solver) de otimização de código aberto para problemas de programação linear (LP), programação inteira mista (MIP) e programação quadrática (QP). Ele é conhecido por sua alta performance e robustez, sendo uma alternativa eficiente para solucionadores comerciais em muitas aplicações. HiGHS pode ser facilmente integrado com ambientes de modelagem como JuMP, proporcionando uma solução completa e de alto desempenho para problemas de otimização. Sua utilização em conjunto com JuMP e SDDP.jl permite a resolução de problemas de programação estocástica multiestágio com eficiência e escalabilidade, tornando-o uma escolha valiosa para este trabalho.